\section{Konsep Dasar Komunikasi Data}

\subsection{Pengertian}
Komunikasi data adalah proses pengiriman dan penerimaan data atau informasi baik menggunakan kabel (wired) maupun tanpa kabel (wireless), dari satu perangkat ke perangkat lain.

\subsection{Tujuan Pengiriman Data}
\begin{enumerate}
\item Memungkinkan pertukaran informasi secara cepat.
\item Mendukung kolaborasi dan berbagi sumber daya.
\item Menjadi dasar bagi berbagai teknologi informasi modern.
\end{enumerate}

\subsection{Komunikasi Data}
$$\text{Pengirim} \to \text{Pesan} \to \text{Penerima}$$

\subsection{Alur Pengiriman Data}
\begin{enumerate}
\item Pembuatan data.
\item Pengubahan data menjadi sinyal.
\item Pengiriman melalui media transmisi.
\item Perubahan sinyal menjadi data.
\end{enumerate}

\subsection{Komponen Dasar Komunikasi Data}
\subsubsection{Pengirim (Sender/Transmitter)}
Pengirim adalah perangkat atau pihak yang mengirim informasi, dapat berupa komputer, smartphone, server, maupun perangkat lainnya yang merupakan sumber asal data.
\subsubsection{Penerima (Receiver)}
Penerima adalah perangkat/pihak yang menerima data yang dikirimkan oleh pengirim yang kemudian akan mengolah dan menampilkan data tersebut agar dapat digunakan oleh pengguna.
\subsubsection{Pesan/Data (Message)}
Pesan/Data adalah informasi yang dikirimkan dari pengirim menuju penerima dapat berupa teks, gambar, audio, video, maupun kombinasi dari berbagai jenis data tersebut.
\subsubsection{Media Transmisi}
Media transmisi adalah jalur fisik atau nirkabel yang dilalui oleh data selama proses pengiriman, misalnya kabel UTP, kabel fiber optik, maupun gelombang radio (wireless) yang menghubungkan pengirim dan penerima.
\subsubsection{Protokol}
Protokol adalah aturan/kesepakatan yang mengatur bagaimana data dikirim, diterima dan di interpretasikan oleh perangkat yang saling berkomunikasi. Tanpa protokol yang sama, perangkat pengirim dan penerima tidak akan saling memahami.

\subsection{Jenis Komunikasi Data}
\subsubsection{Simplex}
Simplex adalah jenis komunikasi data di mana data hanya dapat mengalir dalam satu arah saja dari pengirim menuju penerima. Tanpa adanya jalur balik untuk penerima mengirimkan data kembali ke pengirim. Contohnya adalah siaran radio, televisi.
\subsubsection{Half Duplex}
Half duplex adalah jenis komunikasi data di mana data mengalir dalam dua arah (bolak - balik) namun tidak dapat dilakukan bersamaan dalam mengirim dan menerima data. Contohnya adalah Radio komunikasi (handy talkie/ walkie talkie).
\subsubsection{Full Duplex}
Full duplex adalah jenis komunikasi data di mana data dapat mengalir dalam dua arah secara bersamaan sehingga kedua pihak dapat mengirim dan menerima data pada waktu yang sama tanpa harus bergantian. Contohnya adalah panggilan telepon, panggilan video (video call).
